\section{Part II: From the Mechanism to a Method}
\label{sec:part2-mechanism-to-method}

Part~I leaves the direction as the main loss: centering and a refit of the intercepts handle the threshold, audio needs on the order of $q$ labels per emotion to approach the target's best direction, and a source supplies about ten of them.
Part~II builds the strongest start that sources allow, shows that neither a few labels nor a label-free choice of which utterances to label can move its direction much, and supplies the missing direction information through weak labels on the whole unlabeled pool from a teacher that sees other modalities (Fig.~\ref{fig:method}).

\begin{figure*}[t]
\centering
\FIG{method}{\textwidth}
\caption{The method. Source and target features are centered separately. All sources outside the target corpus form a start $(W_0, b_0)$, from which the student starts with separate $\ell_2$ anchors on the weights and the intercepts. In the basic method a classifier trained on the sources alone (the fused start, the product of the single-modality starts, or their agreement) labels every unlabeled target utterance (for a fused student, also through single-modality students that the fused start teaches) with soft labels whose per-emotion totals are equalized; these weak labels enter the student's loss with total weight $m$ next to the $4K$ true labels. The five values of the start are selected once on pseudo-targets, and the teacher and $m$ for each $K$. The extended version replaces these teachers by the fused student refit on the $K$ labels and the students it teaches, and selects the weighting of the weak labels for each $K$.}
\label{fig:method}
\end{figure*}

\subsection{The Start: Centering and Anchored Transfer}
\label{sec:part2-start}

The student operates on centered features: each source is centered with its own mean, and the target with the mean of its unlabeled pool, which removes most of the threshold loss (Section~\ref{sec:part1-fix}).
The start $(W_0, b_0)$ pools the directions of all sources outside the target corpus.
Each source $s$ contributes the closed-form discriminant direction $W_s(\rho_s) = ((1-\rho_s)\Sigma_s + \rho_s \tau_s I)^{-1}\tilde{M}_s$, where $\tilde{M}_s$ holds its centered class means, $\Sigma_s$ is its covariance, and $\tau_s$ the mean of its eigenvalues; the directions are summed with weights proportional to $N_s^{\gamma}$, and the intercepts assume equal class priors.
The student solves
\begin{equation}
\label{eq:student}
\min_{W, b}\ \sum_{i \in \mathcal{L}} \ell\!\left(W \mathbf{z}_i + b, y_i\right) + \frac{\|W - W_0\|_{M}^2}{2C'} + \frac{\|b - b_0\|^2}{2C_b},
\end{equation}
where $\mathcal{L}$ holds the $4K$ labeled utterances and the norm $\|\cdot\|_M$ uses a metric $M$ built from the scatter of the unlabeled target pool with mixing weight $\rho_M$.
Separate anchors follow Part~I: a few labels set the intercepts well and the weights badly.
The five values $(\gamma, \rho_s, C_b, C', \rho_M)$, three candidates each, are selected on pseudo-targets (Section~\ref{sec:setting-selection}).
For the experiments with weak labels and the comparison, the A, A~(HuBERT), V, and AV students keep $\gamma$ and $\rho_s$ and reselect $(\rho_M, C', C_b)$ on pseudo-targets from a wider grid (Section~S5 of the supplementary material).

Each step of this construction raises UAR: hypothesis transfer from the pooled sources with centering exceeds the best of target-only training, pooling, centered pooling, and pooling with a selected target weight by \NUM{p2-e31-p4-vs-best-audio-lo}--\NUM{p2-e31-p4-vs-best-audio-hi} points for audio at $K \in \{1, 5, 10\}$ (thirty iterations).
Selecting the five values adds \NUM{p2-e39-vs-p4-k5-lo}--\NUM{p2-e39-vs-p4-k5-hi} points at $K=5$ across A, A~(HuBERT), V, and AV, and \NUM{p2-e39-vs-d1-k5-lo}--\NUM{p2-e39-vs-d1-k5-hi} points over a start that weights the sources equally (thirty iterations).
This five-value start is the reference for every addition in Part~II and Section~\ref{sec:comparison}.

\subsection{Where the Direction Cannot Come From}
\label{sec:part2-limits}

Let $W^{*}$ be the student that all labels of the target pool would produce, and let $E_0 = W^{*} - W_0$ be the direction error of the start.
Two facts limit what $K$ labels can do about $E_0$.

\begin{proposition}[Span]
\label{prop:span}
Every row of the solution $W - W_0$ of \eqref{eq:student} lies in the span of the $4K$ labeled points $M^{-1}\mathbf{z}_i$, $i \in \mathcal{L}$.
\end{proposition}

The proposition follows from the representer theorem~\cite{kimeldorf1971results,scholkopf2001generalized} applied to $W - W_0$, since the loss depends on $W$ only through the products $W\mathbf{z}_i$; only the projection of $E_0$ onto a subspace of at most $4K$ dimensions can be corrected, and which part of $E_0$ falls there depends on which utterances carry labels.

\begin{proposition}[Label-free design]
\label{prop:design}
Approximate the loss by a Gaussian likelihood around $W_0$ and the anchor by the prior $W \sim \mathcal{N}(W_0, C' M^{-1})$.
If $E_0$ has no preferred direction under this prior, the expected fraction of $\|E_0\|^2$ that the span of any $4K$ points captures is the same for every choice of points made without their labels.
\end{proposition}

In a Gaussian linear model the posterior covariance does not depend on the observed labels, so neither does the best design~\cite{chaloner1984optimal,chaloner1995bayesian}; a label-free choice can reduce the posterior variance, but without knowledge of the direction of $E_0$ it cannot align the corrected subspace with $E_0$.
Because Proposition~\ref{prop:design} rests on this approximation, we test it directly.
We measured the fraction of $\|E_0\|^2$ captured by the span of $K=5$ labels per emotion under seven label-free selection rules (among them margin, herding, a k-center coreset, and typicality clustering) and under a teaching rule that knows the labels of all candidates.
Relative to random selection, computed as the ratio of the captured fractions averaged over targets, the label-free rules capture \NUM{p2-e58-ratio-unlab-lo}--\NUM{p2-e58-ratio-unlab-hi} times as much, and the teaching rule \NUM{p2-e58-ratio-lll-lo}--\NUM{p2-e58-ratio-lll-hi} times as much (Fig.~S5 of the supplementary material).
Subspaces built without target labels from the disagreement of the sources or from corrections learned on pseudo-targets capture \NUM{p2-e58-r-uaub-lo}--\NUM{p2-e58-r-uaub-hi}\% of $\|E_0\|^2$.
On the five-value student, no label-free rule, including standard active-learning rules, beats random selection by one point at $K=5$ (largest margin \NUM{p2-cond2-e39-maxgain-k5}), whereas replacing a weaker student by the five-value student raises random selection itself by \NUM{p2-cond2-student-random-k5-lo}--\NUM{p2-cond2-student-random-k5-hi} points.
The direction error left by a good start is thus invisible to label-free selection, and its correction needs much more label information than $4K$ labels provide.


\subsection{Weak Labels on the Whole Pool}
\label{sec:part2-weak}

A teacher that predicts every pool utterance supplies label information on thousands of utterances.
Let $Q \in [0,1]^{N_u \times 4}$ hold its class probabilities for the $N_u$ unlabeled pool utterances.
We rescale each column of $Q$ so that every emotion receives the same total weight and scale the matrix to a total weight $m$; the student minimizes \eqref{eq:student} plus $\sum_{i}\sum_{c} Q_{ic}\,\ell(W\mathbf{z}_i + b, c)$ over the pool.
The teacher and $m$, chosen from a small grid that includes no weak labels, are selected for each $K$ on pseudo-targets.

Let $J$ be the UAR of the teacher minus that of the start on the pool, that is, the class-averaged fraction of utterances that the teacher classifies correctly and the start incorrectly, minus the fraction for which the reverse holds.
A student anchored to the start can at best approximately reproduce this advantage, so we expect, without a theorem, that the part of its gain that comes from the teacher's information stays below $J$.
Column normalization adds a label-free effect: the start's own predictions, rebalanced across emotions, already act as a teacher.

Weak labels from a teacher of the student's own modality already improve the student, but such a teacher has a small $J$.
A belief that combines the per-source classifiers, each weighted by its confusion pattern on pseudo-targets, has a $J$ of \NUM{p2-e52-q-minus-p-a} points over the five-value start for WavLM audio.
Its weak labels raise the start by \NUM{p2-e59k-a-k0}, \NUM{p2-e59k-a-k1}, and \NUM{p2-e59k-a-k5} points for audio at $K = 0, 1, 5$ and by \NUM{p2-e59k-v-k0}, \NUM{p2-e59k-v-k1}, and \NUM{p2-e59k-v-k5} points for the face (ten iterations), and the gain vanishes by $K=25$.
For audio at $K=0$, rebalancing the start's own predictions accounts for \NUM{p2-e59k-ctrl-a-k0} points, and the belief's information adds \NUM{p2-e59k-info-a-lo}--\NUM{p2-e59k-info-a-hi} points across $K \le 5$, below its $J$; for text the gain is close to zero.

\subsection{A Teacher from Another Modality}
\label{sec:part2-crossmodal}

When the target utterances come with other modalities, a fused classifier trained on the sources knows a different direction, as in learning with privileged information~\cite{lopezpaz2016unifying,vapnik2009new}.
For a single-modality student, the candidate teachers are the fused start of its configuration, the product of the single-modality starts, and the fused start down-weighted where these starts disagree, all trained on the sources alone (Section~\ref{sec:part2-start}) and never refit; the deployed student uses only its own modality.
In a preliminary check over two iterations in AV, the fused start exceeds the audio start by \NUM{p2-e60pre-juar-av-a-fused} UAR points on the pool, other combinations of modality predictions exceed it by \NUM{p2-e60pre-juar-av-a-lo}--\NUM{p2-e60pre-juar-av-a-hi} points, and the face start has $J = \NUM{p2-e60pre-juar-av-a-v}$ points.
For a fused student, combinations of single-modality starts barely differ from the fused start ($J$ from \NUM{p2-e60pre-juar-av-f-poe} to \NUM{p2-e60pre-juar-av-f-pool} points in the same check), so its candidates add a two-stage reference: single-modality students, trained on the pool and the $K$ labels with the fused start as teacher, are combined, and the product of their predictions either teaches the fused student or replaces it. No teacher of the basic method is refit on the $K$ labels, but the two-stage reference uses them.

At $K \le 5$, cross-modal weak labels raise every student above the five-value start (Fig.~\ref{fig:ladder} shows the audio student).
At $K=0$, the single-modality students gain \NUM{p2-e60-av-a-k0} (audio) and \NUM{p2-e60-av-v-k0} points (face) in AV, \NUM{p2-e60-vt-t-k0} (text) in VT, \NUM{p2-e60-at-t-k0} (text) in AT, and \NUM{p2-e60-avt-a-k0} (audio), \NUM{p2-e60-avt-t-k0} (text), and \NUM{p2-e60-avt-v-k0} points (face) in AVT, and the fused students gain \NUM{p2-e60-f-k0-lo}--\NUM{p2-e60-f-k0-hi} points; all \NUM{p2-e60-m1-n} single-modality students gain at both $K=0$ and $K=5$.
These gains exceed those from a same-modality teacher for audio and text students but not for face students (\NUM{p2-e60-vt-v-k0} and \NUM{p2-e60-avt-v-k0} points in VT and AVT against \NUM{p2-e59k-v-k0} at $K=0$), and in AV the audio student's gain stays below the $J$ of its teacher.
At $K \ge 10$, \NUM{p2-e60-m3-loss-n} of the \NUM{p2-e60-m3-students-n} students fall below the five-value start (worst difference \NUM{p2-e60-m3-worst} points).
The extended family replaces the candidate teachers by the fused student refit on the $K$ labels (the fused start at $K=0$) and a two-stage reference whose single-modality students it teaches, and also selects the weighting of the weak labels (balanced, unbalanced, or hard above a threshold); it removes these losses for the AVT audio and AV fused students evaluated in Section~\ref{sec:comparison} (the VT text and AT audio students were not run with it).

\begin{figure}[t]
\centering
\FIG{ladder}{\columnwidth}
\caption{The method step by step for the audio student. (a)~UAR of target-only training, pooling (mean over single sources), hypothesis transfer from the pooled sources (HT), and the five-value start (seven targets, thirty iterations). (b)~Differences from the five-value start for weak labels from a same-modality belief (seven targets), from the fused teacher of AV (six targets) and AVT (four targets), and from the extended family in AVT (ten iterations); each configuration has its own start.}
\label{fig:ladder}
\end{figure}
